\scanpage{191}
$X$ préschéma connexe, \deleted{\unsure{(réduite, ce qui ne restreint p|(réduction, ce qui ne restreint p|évidemment, ce qui ne restreint p|réduit, ce qui ne restreint pas}}

$1^\circ)$ $X_{\mathrm{réd}}$ de car $p>0$. Alors dans le \unsure{pro-torseur|protorseur} fondamental $(T_\alpha)$, il n'y a pas en général de sections pour les $T_\alpha$. \deleted{\unsure{On}} On dispose uniquement, \deleted{de façon canonique}

(i) des $\xi_{\ell}^{(\alpha)}\in T_\alpha(\mathbf Z_{\ell})$, $l\ne p$, provenant sans doute de $\xi_{\ell}\in T_\alpha(\widetilde{\mathbf Z}_{\ell\mathbf Z})$
\margin{Coh. $l$-adique}
\margin{hensélisé}

\diagram{\begin{tikzcd}[ampersand replacement=\&,column sep=0pt,row sep=3.3em]
\text{(ii) d'un }\xi_p^{(\alpha)}\in T_\alpha( \& \unsure{\mathbb{W}|W}_p(k)\arrow[dr,swap,"{\text{``DeRham-Hodge''}}"] \& )\quad\text{si } \deleted{X\ \unsure{\text{de t.f. sur}|\text{déf.}}}\inserted{\text{$k$ est un}}\ \inserted{\text{corps parfait sur $X$}}\,\deleted{\text{le corps parfait }k} \\
\text{(iii) }\bigl[\,?\quad\text{\uline{réduit}}\ \xi_p^{(\alpha)\,\unsure{0|\circ}}\in T_\alpha( \& \deleted{\unsure{\mathbb{W}}}\ \& k\,)\ \bigr]
\end{tikzcd}}
\margin{Coh. $p$-adique. Structure supplémentaire : opérations \fbox{\uuline{$F$} et \uuline{$V$}} $1/2$ \uuline{linéaires} \uuline{\unsure{reliées|reliant|relativement}} : \unsure{Frobenius uni|Frobenius univ}}
\margin{Coh. Hodge-de Rham Mais ici, il faut \deleted{sans doute} \unsure{donc|faut} aller dans la catégorie $\overline{\unsure{M|\mathcal{M}}}$ dérivée. Tenir compte aussi de la structure \uline{filtrée} en Hodge-De Rham}

[$\xi_p$ provenant sans doute d'un $\xi_p\in T_p(\overset{\sim\!\sim}{\mathbf Z}_{p\mathbf Z})$]\margin{hensélisé strict.}

Si $X=\operatorname{Spec}\mathbf F_q$, il y a en général des \deleted{\illegible{1}} obstructions non triviales dans des $H^2(\mathbf Q,G)$ \ldots

\deleted{\unsure{On}} On utilise des ``foncteurs fibres'' qui induisent sur les motifs effectifs de poids $0$ le foncteur canonique. Alors les \unsure{objets|élts|d pts|points} de $T_\alpha(\mathbf Z_\ell)$ sont sans doute uniques \inserted{à isom près}, (th. de Lang), de $=$ ceux de $T_\alpha(\deleted{\unsure{\mathbf Z}}\mathbb{W}_p(k))$ si $k$ fini (id), \inserted{plus gén.} pour un $k$ parfait de dim \unsure{ch|coh|coh.} $1$ (?). En général, il faut s'attendre que $T_\alpha(k)$ \deleted{\unsure{soit un}} ait plus d'un élément à iso près, mais il est possible que $T_{\deleted{\unsure{\alpha|\infty}}}(k)=\varprojlim T_\alpha(k)$ n'a qu'un élément à iso près, celui donné par la théorie de DeRham-Hodge.
\margin{\unsure{relation} avec les \unsure{problèmes} galoisiens ?? Il doit être \uline{\unsure{équivalent|équivariant}}, cf \uline{Manin}}
\margin{190}
